% Extracto íntegro por residencia del propietario científico definitivo de masas.
\section{Sector nulo y positividad}

La nulidad no se obtiene tomando un límite del carácter positivo. El dominio
se separa primero en
\[
 \mathcal H_{\rm phys}=\mathcal H_0\oplus\mathcal H_+,
\]
con
\[
 \mathcal M|_{\mathcal H_0}=0,
 \qquad
 \mathcal M|_{\mathcal H_+}>0.
\]
El fotón pertenece a la carta nula electromagnética y los gluones a la carta
nula adjunta de color. Sus diferentes representaciones, polarizaciones y
restricciones de calibre permanecen visibles aunque ambas masas sean nulas.

La positividad excluye una especie elemental con \(m^2<0\) dentro del
codominio másico. Un autovalor negativo del Hessiano de un fondo no define un
taquión como partícula: diagnostica la inestabilidad del fondo y exige cambiar
la solución alrededor de la cual se linealiza.

\section{Color, sabores quark y mezcla CKM}

La órbita cromática se realiza en el qutrit de color. Los operadores de Weyl
generan \(\mathfrak{sl}_3(\mathbb C)\); la selección antihermítica de traza
nula produce \(\mathfrak{su}(3)\), y la conexión finita transporta las tres
cartas cromáticas. La masa de sabor debe ser invariante bajo esta acción.

Sea \(U_{\rm col}:SU(3)\to\mathcal U(V_{\rm col})\). La reducción cromática
es la esperanza de Reynolds
\[
 \mathbb E_{\rm col}(T)
 =\int_{SU(3)}U_{\rm col}(g)T U_{\rm col}(g)^*\,dg.
\]
En la representación fundamental irreducible,
\[
 \mathbb E_{\rm col}(\mathcal M_q)
 =m_q(\mu,\mathscr R)I_3.
\]
La masa publicada es, por tanto, independiente de rojo, verde y azul. Las
coordenadas \(\mu\) y \(\mathscr R\) identifican escala y esquema de
renormalización; no son ambigüedades ocultas y deben acompañar siempre a la
cifra.

La dependencia en escala es una sección de un fibrado sobre
\(t=\log\mu\). Si \(\gamma_m(g(t))\) es el endomorfismo inducido por la
conexión cromática, el transporte se define por
\[
 \nabla_{\rm RG}m_q
 =\left(\frac{d}{dt}+\gamma_m(g(t))\right)m_q=0,
\]
\[
 m_q(t_2)
 =\mathcal P\exp\!\left(-\int_{t_1}^{t_2}
 \gamma_m(g(t))\,dt\right)m_q(t_1).
\]
Dos valores pertenecientes a escalas o esquemas distintos no son números
comparables hasta transportarlos a la misma fibra. Esta ecuación cierra el
tipo matemático de la masa corriente: no es un escalar absoluto aislado, sino
una sección covariantemente transportada.

Para mesones y bariones, los proyectores de singlete son
\[
 P_{q\bar q}^{(1)}
 =\frac1{3}|\Omega\rangle\langle\Omega|,
 \qquad
 |\Omega\rangle=\sum_{a=1}^{3}|a\bar a\rangle,
\]
y
\[
 P_{qqq}^{(1)}
 =|\varepsilon\rangle\langle\varepsilon|,
 \qquad
 |\varepsilon\rangle
 =\frac1{\sqrt6}\sum_{a,b,c}\varepsilon_{abc}|abc\rangle.
\]
Ambos eliminan la etiqueta cromática observable sin borrar la órbita interna
que intervino en la construcción.

La mezcla CKM pertenece a un mapa distinto. Si \(B_u\) y \(B_d\) son los
operadores de masa en las bases de tipo arriba y tipo abajo, la matriz unitaria
interna satisface
\[
 B_d^{(u)}=V_{\rm CKM} B_d V_{\rm CKM}^*.
\]
Los ángulos derivados de \(A\), \(C^*\) y \(\Delta_4\), la fase
\(\delta_{\rm CP}\), el invariante de Jarlskog y la identidad del determinante
del conmutador determinan la mezcla sin crear los autovalores de masa. Así se
mantienen separados espectro, color, generación y cambio de base.

\section{Neutrinos, mezcla PMNS y criterio de Majorana}

Los neutrinos son secciones fermiónicas eléctricamente neutras. La
fermionicidad desciende de la holonomía espinorial, la completación exterior y
las relaciones canónicas de anticonmutación; la neutralidad no altera su
estadística. El soporte interno es
\[
 \mathcal H_\nu=
 \bigoplus_{j=1}^{4}\mathcal H_{\nu,G_j},
 \qquad
 \dim\mathcal H_\nu=2+4+6+8=20.
\]
Las cuatro profundidades estructurales no son sabores ni autovalores. El
subespacio físico tridimensional se obtiene mediante el proyector neutral
\(P_\nu\), y el operador de masa es
\[
 M_\nu=P_\nu\mathcal M_\nu P_\nu.
\]
Sus tres autovalores ordenados son las masas; la escala absoluta procede del
mismo reloj másico que fija las restantes ramas, y no de los ángulos de mezcla.

Sea \(M_\ell=P_\ell\mathcal M_\ell P_\ell\) el operador cargado. Denotemos por
\(F_\ell\) y \(F_\nu\) marcos espectrales ortonormales completos, ordenados por
los proyectores espectrales de \(M_\ell\) y \(M_\nu\). La matriz de mezcla es
\[
 U_{\rm PMNS}=F_\ell^*F_\nu.
\]

\begin{theorem}[Cierre espectral de la mezcla leptónica]
Si los dos marcos son completos, \(U_{\rm PMNS}\) es unitaria. Si los
espectros son simples, queda determinada salvo reajustes diagonales de fase y
permutaciones fijadas por el orden espectral. Si existe degeneración, el
resultado exacto es una clase doble
\[
 \prod_\lambda U(m_{\ell,\lambda})
 \backslash U(3) /
 \prod_\mu U(m_{\nu,\mu}),
\]
y no una matriz numérica artificialmente seleccionada.
\end{theorem}

\begin{proof}
La ortonormalidad da
\(U_{\rm PMNS}^*U_{\rm PMNS}=F_\nu^*F_\ell F_\ell^*F_\nu=I_3\).
Los cambios de base dentro de cada subespacio degenerado actúan por la
izquierda o por la derecha, lo que produce exactamente la clase doble
indicada.
\end{proof}

El registro completo determina los elementos de \(M_\nu\) mediante acarreo,
hoja, orientación, vacancia y holonomía. De modo equivalente, puede emplearse
el núcleo hermítico
\[
 K_{\rm reg}(c,d)
 =\sum_{r\in\{\rm car,hoj,ori,vac,hol\}}
 \frac{\langle\Xi_r(c),\Xi_r(d)\rangle}
      {\sum_x\|\Xi_r(x)\|^2},
\]
con omisión de todo sumando de norma nula. Su matriz de Gram es positiva y
construye el operador antes de la diagonalización. La fase de rango uno
\[
 I+(e^{i\pi/6}-1)|b_K\rangle\langle b_K|
\]
permanece como control negativo: su unitariedad formal no basta y su ángulo
\(\theta_{13}\) falsifica su identificación con la mezcla física de rango
completo.

Sea \(G_0\) la matriz tridimensional obtenida al agregar el núcleo anterior
sobre los tres sectores neutrínicos. La condición finita
\[
 \det G_0\ne0
\]
es el certificado exacto de que los observables incluidos en el registro
separan tres direcciones. En tal caso, el factor polar
\[
 U_{G_0}=G_0(G_0^*G_0)^{-1/2}
\]
es unitario y de rango \(3\). Si \(\det G_0=0\), el núcleo declarado no
determina un transporte tridimensional: debe añadirse un observable interno
independiente y volver a calcular la matriz de Gram. No es lícito imponer la
invertibilidad añadiendo un término construido a partir de la matriz de mezcla
que se pretende obtener. La matriz física queda fijada por los marcos
espectrales de \(M_\ell\) y \(M_\nu\); el criterio de Gram certifica suficiencia
del registro, no selecciona por sí solo esos marcos.

Sea \(\mathcal C_\nu\) la autorconjugación antiunitaria del sector neutral. El
criterio de Majorana es
\[
 \mathcal C_\nu^2=I,
 \qquad
 \mathcal C_\nu M_\nu=M_\nu\mathcal C_\nu,
 \qquad
 \mathcal C_\nu H_\nu=H_\nu\mathcal C_\nu,
\]
junto con la existencia de una base real fija en la realificación del
subespacio físico. Cuando se cumplen estas condiciones, cada modo puede
representarse por \(\psi=\mathcal C_\nu\psi\). La rama de Dirac exige, en
cambio, una carga conservada \(Q\) tal que
\[
 \mathcal C_\nu Q\mathcal C_\nu^{-1}=-Q,
 \qquad Q\psi\ne0,
\]
y conserva como modos distintos \(\psi\) y \(\mathcal C_\nu\psi\). Si no se
acredita ninguno de ambos conjuntos de condiciones, el tipo no queda decidido.
La neutralidad, la involución de palabra o el sector de Ising, tomados por
separado, no deciden esta alternativa.

Una profundidad estructural adicional no constituye por sí sola un neutrino
estéril. Un sector estéril candidato queda definido por un proyector no nulo
\(P_s\) que satisface
\[
 P_s P_{\rm act}=0,
 \qquad P_s J_{\text{débil}}=0,
\]
con dinámica definida sobre \(P_s\mathcal H_\nu\). Si además
\([P_s,M_\nu]=0\), el sector es reductor y no mezcla con el sector activo. Si
\(P_sM_\nu P_{\rm act}\ne0\), existe mezcla activa--estéril y necesariamente
\([P_s,M_\nu]\ne0\). Por tanto, la existencia del sector y la existencia de
mezcla son propiedades distintas. Si no existe el proyector débilmente neutro,
el canal grueso adicional es profundidad interna del soporte y no una cuarta
especie.

\section{Bosones masivos y sector escalar}

Los sectores \(W\), \(Z\) y escalar poseen descriptores propios de
representación, ruta y firma. Su realización no exige forzarlos dentro de una
fibra elemental de rango uno. Se construye primero el operador positivo de
fibra y después, en un subespacio \(P_X\mathcal H_X\) de dimensión finita, el
operador efectivo de canales
\[
 H_{X,\rm eff}(z)
 =P_X H_X P_X
 +P_X H_X Q_X(z-Q_X H_X Q_X)^{-1}Q_X H_X P_X.
\]
Si el resolvente reducido admite continuación meromorfa a través del corte del
continuo, los polos son los ceros aislados, contados con multiplicidad, de
\[
 \det\bigl(z-H_{X,\rm eff}(z)\bigr)=0,
 \qquad
 z_X=M_X c_{\rm int}^2-\frac{i}{2}\Gamma_X.
\]
Un canal abierto no garantiza por sí solo la existencia de tales ceros. Cuando
existe un cero con el residuo físico prescrito, la parte real publica la masa de
polo y la parte imaginaria la anchura. Para el
sector escalar, la identificación de una excitación exige además el proyector
escalar y sus acoplamientos de calibre y fermiónicos. Dos polos independientes
determinan dos estados; una segunda etiqueta sin proyector ni polo no crea una
partícula.

\section{Estados compuestos y ley de enlace}

Para constituyentes \(X_1,\ldots,X_n\), la entrada es el registro compuesto
\[
 \mathfrak C_{12}^{(n)}:
 (\Gamma_1,\ldots,\Gamma_n;\mathcal T,\partial)
 \longmapsto\mathcal L_{\rm comp},
\]
no la suma de masas aisladas. La composición binaria
\[
 L_1\star L_2=L_1+L_2+\mathfrak b(L_1,L_2)
\]
es asociativa exactamente cuando
\[
 \mathfrak b(L_1,L_2)+\mathfrak b(L_1\star L_2,L_3)
 =\mathfrak b(L_2,L_3)+\mathfrak b(L_1,L_2\star L_3).
\]
La firma, el refinamiento y el operador compuesto se calculan después del
enlace:
\[
 \mathcal L_{\rm comp}\longmapsto\sigma_{\rm comp}
 \longmapsto\eta_{\rm comp}\longmapsto\mathcal M_{\rm comp}.
\]
El proyector físico reúne singlete de color, espín y paridad, isospín, carga,
simetría de intercambio y canal. La masa estable es un autovalor de la
compresión; cuando se cumplen las hipótesis de continuación meromorfa, la masa
y la anchura de una resonancia proceden de un polo de Feshbach. Con ello
mesones, bariones, estados exóticos, núcleos, átomos y
moléculas poseen el mismo esquema causal sin reducirse a una suma de masas.

Dos árboles de constituyentes representan la misma especie si existe un
unitario que entrelaza simultáneamente sus operadores, proyectores de canal y
polos. En ausencia de tal unitario se conservan como realizaciones distintas,
aunque compartan constituyentes y números cuánticos parciales.

\section{Resonancias, anchuras y supervivencia}

Sea \(S_X\) el proyector sobre el subespacio superviviente y
\(T_X=S_X U S_X\) el operador comprimido de un paso. Si posee un autovalor
dominante simple
\[
 \lambda_X=r_X e^{-i\theta_X},\qquad 0<r_X<1,
\]
fíjese la rama \(\theta_X\in I_\theta\), donde \(I_\theta\) es el intervalo
fundamental de cuasienergía determinado por la carta temporal y su orientación.
Entonces
\[
 M_X=\frac{\hbar_{\rm int}}{t_0c_{\rm int}^2}\theta_X,
 \qquad
 \Gamma_X=-\frac{2\hbar_{\rm int}}{t_0}\log r_X,
 \qquad
 \tau_X=\frac{\hbar_{\rm int}}{\Gamma_X}.
\]
Sin la elección de \(I_\theta\), la fase sólo determina \(\theta_X\) módulo
\(2\pi\) y la masa de cuasienergía no es un escalar bien definido.
La anchura no es una sexta coordenada de la firma másica: procede del módulo
del modo abierto. La igualdad de masa y anchura entre partícula y
antipartícula exige transportar también canales y prescripción causal, no sólo
la paridad de la firma.

\section{Completación multiparticular y estadística}

Sea \(\mathsf{Hol}_{\pm}^{\leq1}\) la categoría monoidal de fibras de una
partícula, con morfismos de refinamiento contractivos, en particular unitarios,
y paridad de holonomía
\(\epsilon\in\{+1,-1\}\). La completación de Fock se define por el funtor
\[
 \mathfrak F(H,\epsilon)=
 \begin{cases}
 \displaystyle\bigoplus_{n\ge0}\operatorname{Sym}^n(H),
     &\epsilon=+1,\\[2mm]
 \displaystyle\bigoplus_{n\ge0}\bigwedge^n\!H,
     &\epsilon=-1,
 \end{cases}
\]
y, para un entrelazador contractivo \(T:H\to H'\),
\[
 \mathfrak F(T)=
 \bigoplus_{n\ge0}\operatorname{Sym}^n(T)
 \quad\text{o}\quad
 \bigoplus_{n\ge0}\bigwedge^n\!T.
\]

\begin{theorem}[Transporte interno de la estadística]
La completación anterior es un funtor de operadores acotados sobre
\(\mathsf{Hol}_{\pm}^{\leq1}\), compatible con refinamiento y
monoidal graduada. En la rama \(\epsilon=-1\) los operadores de creación y
aniquilación satisfacen CAR y los operadores de ocupación cumplen
\(N_s^2=N_s\); en la rama \(\epsilon=+1\) satisfacen CCR. Por tanto, la
exclusión de Pauli y las estadísticas fermiónica y bosónica descienden de la
paridad de holonomía una vez fijada la representación de una partícula.
\end{theorem}

\begin{proof}
Las potencias simétricas y exteriores preservan composiciones e identidades;
la contractividad de \(T\) garantiza que su suma directa define un operador
acotado. Para morfismos no contractivos, la segunda cuantización requiere un
dominio denso explícito y no pertenece a esta categoría de operadores acotados.
La regla de signos del producto exterior produce CAR y la idempotencia de la
ocupación; el cociente simétrico produce CCR. La segunda cuantización de los
morfismos de refinamiento entrelaza los correspondientes operadores.
\end{proof}

Esta demostración cierra el transporte algebraico interno. La identificación
con el teorema relativista de espín--estadística requiere, además, el puente
lorentziano y la localidad declarados en Mecánica Dimensional; no deben
confundirse ambos enunciados.

Para un Hamiltoniano comprimido \(H_{X,V}\) sobre un volumen o refinamiento
finito \(V\), el estado térmico está completamente definido por
\[
 Z_{X,V}(\beta)=\operatorname{Tr}e^{-\beta H_{X,V}},
 \qquad
 \rho_{X,V,\beta}=Z_{X,V}(\beta)^{-1}e^{-\beta H_{X,V}}.
\]
El límite termodinámico existe cuando la presión
\[
 p_X(\beta)=\lim_{V\nearrow\infty}
 \frac1{\beta |V|}\log Z_{X,V}(\beta)
\]
existe y es independiente de la sucesión cofinal. La condensación bosónica se
decide por la saturación de la densidad de estados excitados; no se proclama a
partir de la mera existencia de una completación simétrica. Así, el problema
queda expresado mediante una condición matemática verificable para el
Hamiltoniano y el límite, en lugar de una atribución nominal.

\section{Sectores topológicos y estados virtuales}

En un sector topológico el codominio primario es una categoría trenzada, no
una recta de masas. Para objetos simples \(a,b,c\), los datos físicos incluyen
\[
 N_{ab}^{c},\qquad F^{abc}_{d},\qquad R^{ab}_{c},
\]
con ecuaciones pentagonal y hexagonal. Una masa escalar existe únicamente
cuando una realización Hamiltoniana proporciona un proyector espectral y una
brecha. Los sectores Fibonacci, Ising/Majorana y \(\mathbb Z_6\) quedan así
cerrados como codominios categóricos; no se les asigna una cifra universal
por analogía con una partícula elemental.

En el sector \(\mathbb Z_6\), la fusión y el trenzado son
\[
 a\otimes b=a+b\pmod 6,
 \qquad
 R_{ab}=\zeta_6^{ab},
 \qquad \zeta_6=e^{2\pi i/6}.
\]
Para Fibonacci,
\[
 \tau\otimes\tau=\mathbf 1\oplus\tau,
 \qquad
 F_\tau^{\tau\tau\tau}=
 \begin{pmatrix}
 \varphi^{-1}&\varphi^{-1/2}\\
 \varphi^{-1/2}&-\varphi^{-1}
 \end{pmatrix},
\]
\[
 R_{\mathbf1}^{\tau\tau}=e^{-4\pi i/5},
 \qquad
 R_\tau^{\tau\tau}=e^{3\pi i/5}.
\]
El sector de Ising posee objetos \(\mathbf1,\psi,\sigma\) y reglas
\[
 \psi\otimes\psi=\mathbf1,
 \qquad
 \psi\otimes\sigma=\sigma,
 \qquad
 \sigma\otimes\sigma=\mathbf1\oplus\psi,
\]
\[
 F_\sigma^{\sigma\sigma\sigma}
 =\frac1{\sqrt2}
 \begin{pmatrix}1&1\\1&-1\end{pmatrix},
 \qquad
 R_{\mathbf1}^{\sigma\sigma}=e^{-\pi i/8},
 \qquad
 R_\psi^{\sigma\sigma}=e^{3\pi i/8}.
\]
Estas matrices satisfacen las relaciones de coherencia en las convenciones
indicadas. La duplicación
\(\mathcal C\boxtimes\mathcal C^{\rm rev}\) conserva ambas orientaciones y
evita introducir una quiralidad por elección externa.

Un estado virtual es una sección de persistencia finita en una amplitud o un
resolvente. No pertenece al espectro asintótico y no exige una masa de polo
real. Sus parámetros son la virtualidad, el canal y la duración de
persistencia. Esta definición evita convertir una línea interna de cálculo en
una nueva especie.

\Needspace{25\baselineskip}
\section{Gravitación, torsión y modo de espín 2}

La gravitación entra en Mecánica Dimensional mediante conexión, curvatura y
torsión. No requiere una partícula elemental de espín \(2\) como dato de
partida. Sea \(\mathcal S_{\rm grav}\) la acción geométrica y
\(\mathcal H_{\rm TT}\) el subespacio transversal y sin traza de su Hessiano
en una solución. El observable linealizado es
\[
 \mathcal K_{\rm TT}
 =P_{\rm TT}\,\delta^2\mathcal S_{\rm grav}\,P_{\rm TT}.
\]
Un modo colectivo de espín \(2\) existe como partícula asintótica únicamente
si el resolvente de \(\mathcal K_{\rm TT}\) posee un polo físico con residuo
positivo y respeta las restricciones. Si no existe tal polo, la interacción
gravitatoria permanece íntegramente en la conexión y la torsión. Así, el
gravitón no es una pieza necesaria de la ley de masas, y su eventual aparición
queda decidida por un criterio espectral, no por una casilla nominal.
